\section{Beyond spiders: rooted trees and bounded cyclic blocks}
\label{sec:tree-and-block-results}

The spider calculation suggests two possible notions of generalization.

\begin{enumerate}[leftmargin=*]
\item \emph{SOR generalization} asks whether the same signed iterative
      recurrence has accelerated semantic convergence and locally maintainable
      support.
\item \emph{target-problem generalization} asks whether some fully charged
      local algorithm reaches the desired PPR work scale on the larger graph
      class.
\end{enumerate}

The first remains incomplete beyond a fixed face.  The second is already
proved for substantially more than spiders through persistent responses.

\subsection{Arbitrary rooted trees}

For a finite tree rooted at its single seed, the proof-owning
\texttt{propagate\_settle\_framework} note constructs lazy exact subtree
responses.  Conditioned on its parent value, every subtree sends a monotone
piecewise-affine scalar response.  Child breakpoints are merged lazily, so an
active label is propagated only along its root path; one final top-down solve
materializes the solution.

Let $S^\rho=\operatorname{supp}(x^\rho)$ and
$R_\star=\max_{v\in S^\rho}\operatorname{dist}(o,v)$.  The imported results
give
\begin{align}
 R_\star
 &\leq1+
 \frac{\log_+((1-\alpha)/(\alpha\rho))}
      {\log((1+\sqrt\alpha)/(1-\sqrt\alpha))},
 \label{eq:tree-support-radius-import}\\
 W_{\rm tree}
 &\leq O\!\left(
  \log(2+\Delta_\star)
  \left[
   \operatorname{cvol}(S^\rho)
   +\sum_{v\in S^\rho}\operatorname{dist}(o,v)
  \right]
 \right)
 \notag\\
 &=\widetilde O\!\left(
  \frac{1+\operatorname{cvol}(S^\rho)}{\sqrt\alpha}
 \right).
 \label{eq:tree-response-import}
\end{align}
The algorithm scans active adjacency lists, emits every active label, charges
response-event routing, performs final recovery, and checks the exterior KKT
condition.  It does not require the support or radius as advice.

\begin{theorem}[Semantic PPR at the product scale on rooted trees]
\label{thm:tree-semantic-ppr}
On every finite unweighted tree with one seed, exact arithmetic and ordinary
adjacency-list access suffice to return a sparse $\widehat x$ with
\begin{equation}
 \|\widehat x-x^0\|_{D,\infty}
 \leq\varepsilon_{\rm ppr}
 \label{eq:tree-semantic-target}
\end{equation}
in
\begin{equation}
 W_{\rm tree}
 =\widetilde O\!\left(
  \frac1{\sqrt\alpha\,\varepsilon_{\rm ppr}}
 \right)
 \label{eq:tree-semantic-work}
\end{equation}
charged work, including discovery, exact solution, certification, and output.
\end{theorem}

\begin{proof}
Set $\rho=\varepsilon_{\rm ppr}$ and return the exact RPPR point
$\widehat x=x^\rho$ from the imported response algorithm.  By
\Cref{thm:rppr-ppr-bias}, its semantic PPR error is at most $\rho$.  The
source support bound gives
$\operatorname{vol}(S^\rho)\leq1/\rho$.  If the support is nonempty, every
degree is at least one and hence
\begin{equation*}
 \operatorname{cvol}(S^\rho)
 =|S^\rho|+\operatorname{vol}(S^\rho)
 \leq2\operatorname{vol}(S^\rho)
 \leq\frac2\rho.
\end{equation*}
The empty-support case is constant work.  Substitute into
\eqref{eq:tree-response-import}; its soft factors come from the displayed
support radius and local heaps.  Output is already charged by the imported
theorem.
\end{proof}

This is an end-to-end theorem for a broad infinite graph class.  It is not an
SOR theorem: the response representation retains exact piecewise-affine
subtree maps.  On a known fixed tree face, \Cref{cor:fixed-face-semantic-sor}
still gives an accelerated iterative tail, but running that tail after the
response algorithm has already computed the exact solution would be
redundant.

\subsection{Graphs assembled from bounded cyclic blocks}

Trees cease to be the right boundary when small cycles are attached at
articulation vertices.  Count each bridge as a two-vertex block and suppose
every biconnected block in the seed component has at most $q$ vertices.  The
proof-owning \texttt{delayed\_reflection\_ladder} note maintains the
block--cut forest under ordinary adjacency exposure.  When a new edge merges
blocks, the newly activated coordinates are still zero at that event, so
past response pieces remain valid; only future pieces in the merged block are
rebuilt.

The imported representation-free theorem solves exact single-seed RPPR in
\begin{equation}
 W_{\rm block}
 =\widetilde O\!\left(
  \frac{q^3}{\rho\sqrt\alpha}
 \right)
 \label{eq:bounded-block-rppr-import}
\end{equation}
degree work and $O(q\operatorname{vol}(S^\rho))$ storage.  It requires no
supplied block decomposition, support, radius, or confinement oracle.

\begin{corollary}[Semantic PPR on bounded-block graphs]
\label{cor:bounded-block-semantic-ppr}
Under the bounded-block hypothesis above, choosing
$\rho=\varepsilon_{\rm ppr}$ returns semantic PPR output in
\begin{equation}
 W_{\rm block}
 =\widetilde O\!\left(
  \frac{q^3}{\sqrt\alpha\,\varepsilon_{\rm ppr}}
 \right).
 \label{eq:bounded-block-ppr-work}
\end{equation}
In particular, constant-size blocks meet the desired product scale, and
$q=\log^{O(1)}(1/(\alpha\varepsilon_{\rm ppr}))$ changes only soft factors.
\end{corollary}

\begin{proof}
Apply \Cref{thm:rppr-ppr-bias} to the exact RPPR output and substitute
$\rho=\varepsilon_{\rm ppr}$ into
\eqref{eq:bounded-block-rppr-import}.
\end{proof}

\subsection{The graph-family ladder now established}

The logical status is summarized below.

\begin{center}
\begin{tabular}{p{0.16\textwidth}p{0.22\textwidth}p{0.25\textwidth}p{0.19\textwidth}}
\toprule
Graph class & Strongest proved primitive & Charged semantic PPR scale & Main
qualification \\
\midrule
Center-star & Signed coordinate SOR &
$\Theta(1/(\sqrt\alpha\varepsilon_{\rm ppr}))$ in the declared class &
Promised center-star; exact semantic stop \\
Fixed bipartite face & Red--black SOR &
$\widetilde O(\operatorname{vol}(U)/\sqrt\alpha)$ & Face and exterior test
supplied \\
Hub-rooted spider & Kinetic scalar response &
$O(\varepsilon_{\rm ppr}^{-1}\log(2+B))$ & Exact arithmetic; tree geometry
\\
Arbitrary rooted tree & Lazy subtree response &
$\widetilde O(1/(\sqrt\alpha\varepsilon_{\rm ppr}))$ & Exact arithmetic;
single seed \\
Blocks of size $q$ & Dynamic block response &
$\widetilde O(q^3/(\sqrt\alpha\varepsilon_{\rm ppr}))$ & Every biconnected
block has size at most $q$ \\
Arbitrary graph & None yet & Target remains
$\widetilde O(1/(\sqrt\alpha\varepsilon_{\rm ppr}))$ & Large cyclic cores
remain open \\
\bottomrule
\end{tabular}
\end{center}

This ladder changes the research priority.  Proving another tree-only SOR
rate would refine the iterative story, but it would not enlarge the graph
class on which the target problem is solved.  The next genuinely broader
step must address a large biconnected core.
